\begin{abstract}
Semi-supervised learning (SSL) methods for 3D medical image segmentation discard the unlabeled training data at convergence, even though the converged model's predictions on it still contain residual uncertainty concentrated at organ boundaries. We measure this residual entropy directly on the publicly released BCP checkpoints for Pancreas-CT 20\%, LA 5\%, and LA 10\% and find that \textbf{$89$--$95\%$ of the network's total prediction entropy is carried by a thin uncertain-voxel shell occupying only $5$--$9\%$ of each volume}, whereas a supervised-only baseline trained on the same labels carries only $\sim\!55\%$ of its entropy in the corresponding shell. We introduce \textbf{Post-hoc Entropy Minimization (PEM)}: a label-free, pseudo-label-free, architecture-free fine-tuning step that takes any converged SSL checkpoint and minimizes prediction entropy on the unlabeled training set for $2$--$5$ epochs. Under a fixed-budget protocol with no test-set hyperparameter tuning and three random seeds per configuration, PEM improves Dice by $+0.6$ to $+1.4$ and HD95 by $27$--$43\%$ in under $5.2$ minutes per run on a single GPU; all three gains pass a per-case Wilcoxon signed-rank test at $p<0.025$. We give a theoretical account of why post-hoc entropy minimization is well-conditioned on a converged checkpoint and ill-conditioned elsewhere, via an implicit boundary-shell gating argument based on the $\partial H/\partial z=-z\,p(1{-}p)$ gradient of the binary entropy, and we define a scalar, label-free diagnostic $\rho_f(\tau)$ that predicts applicability from a single forward pass. Against three controlled post-hoc baselines (temperature scaling, pseudo-label fine-tune, self-consistency), PEM attains the largest Dice gain on the two less-saturated configurations, and pseudo-label fine-tuning catches up in the saturation limit in a manner quantitatively predicted by the theory.
\end{abstract}
